%% ma: L12-B15-0023
%% lop: 12
%% chuong: 5
%% bai: 15
%% dang: 12.15.9
%% loai: TL
%% muc: VD
%% dap_an: "-"
%% nguon: "(NBV)-10 ĐỀ KIỂM TRA HK2-TOÁN 12-DUNG CHUNG 3 SÁCH.pdf (H03, Đề số 3), Phần 4 Câu 2 (tự luận)"
%% kien_thuc: "Bài 14 (phương trình mặt phẳng, điểm thuộc mặt phẳng), Bài 15 (phương trình đường thẳng, điểm thuộc đường thẳng), tọa độ trung điểm (Bài 7)"
%% trang_thai: chinh-thuc
%% ghi_chu: "Đáp án gốc khớp về phương trình đường thẳng, nhưng lời giải gốc ghi nhầm tên điểm (M(1;0;-1), N(3;2;5)); đúng là M(3;2;5), N(-1;4;-1)."
\begin{ex}
Trong không gian $Oxyz$, cho mặt phẳng $(P)\colon x+y+2z-15=0$, điểm $A(1;3;2)$ và đường thẳng $d\colon\begin{cases}x=1+t\\y=2-t\\z=3+2t\end{cases}$. Tìm phương trình đường thẳng $\Delta$ cắt $(P)$ và $d$ lần lượt tại $M$, $N$ sao cho $A$ là trung điểm của đoạn $MN$.
\loigiai{$N\in d\Rightarrow N(1+t;2-t;3+2t)$. Vì $A$ là trung điểm của $MN$ nên $M=2A-N=(1-t;4+t;1-2t)$.

$M\in(P)$: $(1-t)+(4+t)+2(1-2t)-15=0\Leftrightarrow-8-4t=0\Leftrightarrow t=-2$. Suy ra $N(-1;4;-1)$, $M(3;2;5)$.

$\overrightarrow{MN}=(-4;2;-6)$ nên $\Delta$ đi qua $M(3;2;5)$ với vectơ chỉ phương $\vec a=(2;-1;3)$:
\[\Delta\colon\dfrac{x-3}{2}=\dfrac{y-2}{-1}=\dfrac{z-5}{3}.\]}
\end{ex}
